% ═══════════════════════════════════════════════════════════════════════════
%  فصل ۳ — الگوهای تحلیلی: موفقیت، شکست و ارزش‌سنجی
% ═══════════════════════════════════════════════════════════════════════════
\chapter{الگوهای تحلیلی: موفقیت، شکست و ارزش‌سنجی}
\label{ch:models}

\begin{tcolorbox}[mybox, title={چکیده‌ی فصل}]
این فصل الگوهای نظری و تحلیلی برای ارزیابی مداخلات خارجی و نجات‌طلبی را معرفی و نقد می‌کند. از الگوی واقع‌گرایانه‌ی «هزینه-فایده» تا الگوهای نهادگرایانه، سازه‌انگارانه و پسااستعماری، و در نهایت الگوی ترکیبی پیشنهادی این پژوهش ارائه می‌شود.
\end{tcolorbox}

% ─────────────────────────────────────────────────────────────────────────
\section{چرا به الگوی تحلیلی نیاز داریم؟}
\label{sec:ch03-why}
% ─────────────────────────────────────────────────────────────────────────

مطالعه‌ی موردی بدون چارچوب تحلیلی، انبوهه‌ای از واقعیت‌های پراکنده است. \thinker{الکساندر جرج} (\lr{Alexander George}) و \thinker{اندرو بنت} (\lr{Andrew Bennett}) تأکید می‌کنند که «مطالعه‌ی موردی ساختاریافته و متمرکز» (\lr{structured focused comparison}) مستلزم سؤالات استاندارد و معیارهای سنجش یکسان برای همه‌ی موارد است.%
\footnote{\lr{George, Alexander L., and Andrew Bennett. \textit{Case Studies and Theory Development in the Social Sciences}. MIT Press, 2005, pp.~67--72.}}

الگوی تحلیلی ما باید بتواند:
\begin{enumerate}[nosep, label=\textbf{\arabic*.}]
  \item \textbf{توصیف} کند: چه اتفاقی افتاد؟
  \item \textbf{تبیین} کند: چرا اتفاق افتاد؟
  \item \textbf{ارزیابی} کند: نتیجه چه بود؟
  \item \textbf{پیش‌بینی} کند: در شرایط مشابه چه خواهد شد؟
\end{enumerate}

% ─────────────────────────────────────────────────────────────────────────
\section{الگوی ۱: تحلیل واقع‌گرایانه‌ی هزینه-فایده}
\label{sec:ch03-realist}
% ─────────────────────────────────────────────────────────────────────────

\subsection{مبانی نظری}
\label{subsec:ch03-realist-theory}

الگوی واقع‌گرایانه (\lr{Realist Cost-Benefit Model}) بر پیش‌فرض‌های زیر استوار است:

\begin{itemize}[nosep, rightmargin=1em]
  \item دولت‌ها بازیگران عقلانی هستند که منافع ملی خود را حداکثر می‌کنند.
  \item مداخله‌ی خارجی زمانی رخ می‌دهد که هزینه‌های آن (نظامی، اقتصادی، اعتباری) کمتر از فوایدش (ژئوپولیتیکی، اقتصادی، امنیتی) باشد.
  \item دعوت‌کنندگان داخلی زمانی به بیگانه متوسل می‌شوند که هزینه‌ی «عدم اقدام» از هزینه‌ی «وابستگی» بیشتر باشد.
\end{itemize}

\thinker{پاتریک ریگان} (\lr{Patrick Regan}) در پژوهش آماری جامع خود روی ۱۵۰ جنگ داخلی (۱۹۴۴–۱۹۹۹) نشان داده که مداخله‌ی خارجی به‌طور متوسط \textbf{مدت جنگ‌های داخلی را کاهش نمی‌دهد} بلکه در بسیاری موارد آن را طولانی‌تر می‌کند.%
\footnote{\lr{Regan, Patrick M. "Third-Party Interventions and the Duration of Intrastate Conflicts." \textit{Journal of Conflict Resolution} 46, no.~1 (2002): 55--73.}}

\subsection{ماتریس هزینه-فایده}
\label{subsec:ch03-realist-matrix}

\begin{table}[htbp]
\centering
\caption{ماتریس هزینه-فایده‌ی مداخله از دید بازیگران مختلف}\label{tab:ch03-cost-benefit}
\begin{adjustbox}{max width=\linewidth}
\begin{tabular}{
  >{\raggedleft\arraybackslash\bfseries}p{2.5cm}
  >{\raggedleft\arraybackslash}p{3.5cm}
  >{\raggedleft\arraybackslash}p{3.5cm}
  >{\raggedleft\arraybackslash}p{3.5cm}}
\toprule
\rowcolor{PrimaryDeep}
\textcolor{white}{بُعد} &
\textcolor{white}{دعوت‌کننده} &
\textcolor{white}{مداخله‌گر} &
\textcolor{white}{حکومت هدف} \\
\midrule

\multicolumn{4}{r}{\cellcolor{AccentGreen!15}\textbf{\color{AccentGreen}فواید بالقوه}} \\
\midrule
نظامی &
سقوط رژیم ستم‌گر &
پایگاه‌های نظامی، نفوذ &
— \\
\rowcolor{NeutralLight}
سیاسی &
قدرت‌یابی &
رژیم دست‌نشانده &
— \\
اقتصادی &
بازسازی (احتمالی) &
امتیازات اقتصادی، منابع &
— \\
\rowcolor{NeutralLight}
اعتباری &
مشروعیت (کوتاه‌مدت) &
وجهه‌ی «ناجی» &
— \\

\midrule
\multicolumn{4}{r}{\cellcolor{AccentRed!15}\textbf{\color{AccentRed}هزینه‌های بالقوه}} \\
\midrule
نظامی &
تخریب، تلفات غیرنظامی &
تلفات نظامی، هزینه‌ی لجستیک &
شکست نظامی \\
\rowcolor{NeutralLight}
سیاسی &
وابستگی، از دست‌دادن حاکمیت &
مقاومت داخلی، باتلاق &
سقوط \\
اقتصادی &
ویرانی زیرساخت &
هزینه‌ی نجومی بازسازی &
از دست‌دادن منابع \\
\rowcolor{NeutralLight}
اعتباری &
انگ «خیانت» و «وابستگی» &
انگ «امپریالیسم» &
انگ «استبداد» \\

\bottomrule
\end{tabular}
\end{adjustbox}
\end{table}

\subsection{نقد الگوی واقع‌گرایانه}
\label{subsec:ch03-realist-critique}

\begin{itemize}[nosep, rightmargin=1em]
  \item فرض «عقلانیت» بازیگران اغلب غیرواقعی است. تصمیم‌ها تحت فشار، هیجان و اطلاعات ناقص گرفته می‌شوند.%
  \footnote{\lr{Jervis, Robert. \textit{Perception and Misperception in International Politics}. Princeton UP, 1976.}}
  \item هزینه‌ها و فواید «غیرمادی» (هویت، شأن، عدالت) در این الگو نادیده گرفته می‌شوند.
  \item تمرکز بر دولت‌ها به‌عنوان تنها بازیگر، نقش جنبش‌های اجتماعی و شبکه‌های فراملی را نادیده می‌گیرد.
\end{itemize}

% ─────────────────────────────────────────────────────────────────────────
\section{الگوی ۲: تحلیل نهادگرایانه}
\label{sec:ch03-institutional}
% ─────────────────────────────────────────────────────────────────────────

الگوی نهادگرایانه (\lr{Institutional Analysis}) بر نقش نهادهای بین‌المللی (سازمان ملل، ناتو، اتحادیه‌ی آفریقایی، و غیره) در تسهیل، مشروعیت‌بخشی یا مهار مداخلات تمرکز دارد.

\thinker{رابرت کیوهین} (\lr{Robert Keohane}) استدلال می‌کند که نهادهای بین‌المللی می‌توانند هزینه‌ی مبادلات را کاهش دهند، اطلاعات را افزایش دهند و تعهدات را پایدار کنند.%
\footnote{\lr{Keohane, Robert O. \textit{After Hegemony}. Princeton UP, 1984.}}

\subsection{شاخص‌های نهادی موفقیت مداخله}
\label{subsec:ch03-inst-indicators}

\begin{tcolorbox}[defbox, title={شاخص‌های کلیدی}]
\begin{enumerate}[nosep, label=\textbf{\arabic*.}]
  \item \textbf{مجوز نهادی:} آیا مداخله مجوز شورای امنیت یا سازمان منطقه‌ای داشته؟
  \item \textbf{چندجانبه‌گرایی:} آیا چندین دولت مشارکت داشته‌اند؟
  \item \textbf{نظارت و پاسخ‌گویی:} آیا سازوکار نظارت بر رفتار مداخله‌گر وجود داشته؟
  \item \textbf{برنامه‌ی خروج:} آیا استراتژی خروج از پیش تعیین شده بود؟
  \item \textbf{بازسازی نهادی:} آیا پس از مداخله، نهادسازی صورت گرفته؟
\end{enumerate}
\end{tcolorbox}

% ─────────────────────────────────────────────────────────────────────────
\section{الگوی ۳: تحلیل سازه‌انگارانه}
\label{sec:ch03-constructivist}
% ─────────────────────────────────────────────────────────────────────────

الگوی سازه‌انگارانه (\lr{Constructivist Model}) بر نقش \concept{هنجارها}، \concept{هویت‌ها} و \concept{روایت‌ها} در شکل‌دهی به مداخلات تأکید دارد.

\thinker{مارتا فینه‌مور} نشان داده که «انگیزه‌های» مداخله در طول زمان تحول یافته‌اند — نه به‌خاطر تغییر در منافع مادی، بلکه به‌خاطر تغییر در هنجارهای بین‌المللی:

\begin{longtable}{
  >{\raggedleft\arraybackslash\bfseries}p{3cm}
  >{\raggedleft\arraybackslash}p{3.5cm}
  >{\raggedleft\arraybackslash}p{3.5cm}
  >{\raggedleft\arraybackslash}p{3cm}}
\caption{تحول هنجاری مداخلات (بر اساس فینه‌مور)}\label{tab:ch03-finnemore}\\
\toprule
\rowcolor{PrimaryDeep}
\textcolor{white}{\textbf{دوره}} &
\textcolor{white}{\textbf{هنجار غالب}} &
\textcolor{white}{\textbf{بهانه‌ی مداخله}} &
\textcolor{white}{\textbf{نمونه}} \\
\midrule
\endfirsthead
\toprule
\rowcolor{PrimaryDeep}
\textcolor{white}{\textbf{دوره}} &
\textcolor{white}{\textbf{هنجار غالب}} &
\textcolor{white}{\textbf{بهانه‌ی مداخله}} &
\textcolor{white}{\textbf{نمونه}} \\
\midrule
\endhead
\bottomrule
\endlastfoot

قرن ۱۹ &
حمایت از هم‌کیشان مسیحی &
کشتار مسیحیان &
مداخله در عثمانی%
\footnote{\lr{Finnemore, Martha. \textit{The Purpose of Intervention}. Cornell UP, 2003, Ch.~2.}} \\
\rowcolor{NeutralLight}
دوران جنگ سرد &
جلوگیری از گسترش کمونیسم/سرمایه‌داری &
تهدید ایدئولوژیک &
کره، ویتنام \\
پس از جنگ سرد &
حمایت از حقوق بشر جهان‌شمول &
نسل‌کشی، پاک‌سازی قومی &
سومالی، کوزوو \\
\rowcolor{NeutralLight}
پس از ۲۰۰۱ &
مبارزه با تروریسم + دموکراسی‌سازی &
تهدید تروریستی &
افغانستان، عراق \\
پس از ۲۰۱۱ &
بحران هنجاری: شکست \lr{R2P} در سوریه &
بی‌اعتمادی فزاینده &
سوریه، یمن \\

\end{longtable}

% ─────────────────────────────────────────────────────────────────────────
\section{الگوی ۴: تحلیل پسااستعماری}
\label{sec:ch03-postcolonial}
% ─────────────────────────────────────────────────────────────────────────

الگوی پسااستعماری (\lr{Postcolonial Analysis}) نجات‌طلبی و مداخله را در بستر \concept{ساختارهای قدرت نابرابر جهانی} تحلیل می‌کند.

\subsection{مفاهیم کلیدی}
\label{subsec:ch03-postcolonial-concepts}

\begin{tcolorbox}[defbox, title={مفاهیم کلیدی الگوی پسااستعماری}]
\begin{description}[nosep, style=nextline]
  \item[\concept{شرق‌شناسی} (\lr{Orientalism})]
  بازنمایی «شرق» به‌عنوان دیگری ناتوان و نیازمند نجات — مشروعیت‌بخش مداخله (\thinker{سعید}).
 
  \item[\concept{درونی‌سازی سلطه} (\lr{Internalized Oppression})]
  قربانی ستم، خود را ناتوان می‌بیند و «ناجی» را در بیگانه جست‌وجو می‌کند (\thinker{فانون}).
  \item[\concept{وابستگی ساختاری} (\lr{Structural Dependency})]
  مداخله، ساختار اقتصادی-سیاسی وابسته ایجاد می‌کند که حتی پس از خروج مداخله‌گر باقی می‌ماند (\thinker{آندره گوندر فرانک}).%
  \footnote{\lr{Frank, Andre Gunder. \textit{Capitalism and Underdevelopment in Latin America}. Monthly Review Press, 1967.}}
  \item[\concept{خشونت معرفتی} (\lr{Epistemic Violence})]
  مداخله‌گر نه‌فقط نظامی بلکه فرهنگی-معرفتی نیز عمل می‌کند: تعریف «پیشرفت»، «دموکراسی» و «حقوق بشر» را تحمیل می‌کند (\thinker{گایاتری اسپیواک}).%
  \footnote{\lr{Spivak, Gayatri Chakravorty. "Can the Subaltern Speak?" In \textit{Marxism and the Interpretation of Culture}, ed. Nelson and Grossberg. U of Illinois P, 1988, pp.~271--313.}}
\end{description}
\end{tcolorbox}

\subsection{نقد پسااستعماری مداخله‌ی انسان‌دوستانه}
\label{subsec:ch03-postcolonial-critique}

\thinker{سواتی پاراشار} (\lr{Swati Parashar}) استدلال می‌کند که روایت «نجات زنان» از دست رژیم‌های ستم‌گر — مانند روایت امریکا در افغانستان — بازتولید الگوی استعماری «مرد سفید زن رنگین‌پوست را از مرد رنگین‌پوست نجات می‌دهد» است.%
\footnote{\lr{Parashar, Swati. "Feminist International Relations and Women Militants." \textit{Cambridge Review of International Affairs} 22, no.~2 (2009): 235--256.}}

\thinker{تالال اسد} (\lr{Talal Asad}) نیز نشان می‌دهد که مفهوم «جنگ عادلانه» در سنت غربی، ذاتاً بازتاب سلسله‌مراتب قدرت جهانی است و نمی‌تواند معیاری جهان‌شمول برای ارزیابی مداخلات باشد.%
\footnote{\lr{Asad, Talal. \textit{On Suicide Bombing}. Columbia UP, 2007, pp.~15--62.}}

% ─────────────────────────────────────────────────────────────────────────
\section{الگوی ۵: تحلیل روان‌شناسی سیاسی}
\label{sec:ch03-psychological}
% ─────────────────────────────────────────────────────────────────────────

الگوی روان‌شناسی سیاسی بر فرایندهای شناختی و عاطفی تصمیم‌گیران و توده‌ها تمرکز دارد:

\subsection{سوگیری‌های شناختی در نجات‌طلبی}
\label{subsec:ch03-biases}

\begin{longtable}{
  >{\raggedleft\arraybackslash\bfseries}p{3cm}
  >{\raggedleft\arraybackslash}p{4cm}
  >{\raggedleft\arraybackslash}p{5.5cm}}
\caption{سوگیری‌های شناختی مؤثر در نجات‌طلبی}\label{tab:ch03-biases}\\
\toprule
\rowcolor{PrimaryDeep}
\textcolor{white}{\textbf{سوگیری}} &
\textcolor{white}{\textbf{تعریف}} &
\textcolor{white}{\textbf{نمود در نجات‌طلبی}} \\
\midrule
\endfirsthead
\toprule
\rowcolor{PrimaryDeep}
\textcolor{white}{\textbf{سوگیری}} &
\textcolor{white}{\textbf{تعریف}} &
\textcolor{white}{\textbf{نمود در نجات‌طلبی}} \\
\midrule
\endhead
\bottomrule
\endlastfoot

آرزواندیشی (\lr{Wishful Thinking}) &
باور به نتیجه‌ی مطلوب بدون شواهد کافی &
«امریکا/روسیه/... ما را نجات خواهد داد و سپس خواهد رفت»%
\footnote{\lr{Jervis, Robert. \textit{Perception and Misperception in International Politics}. Princeton UP, 1976, Ch.~4.}} \\
\rowcolor{NeutralLight}
خام‌اندیشی (\lr{Naive Realism}) &
فرض بر اینکه دیگران جهان را همان‌گونه می‌بینند &
«مداخله‌گر منافع ما را همان‌قدر که ما اهمیت می‌دهیم، اهمیت خواهد داد» \\
اثر هاله‌ای (\lr{Halo Effect}) &
تعمیم یک ویژگی مثبت به کل شخصیت &
«چون دموکراسی دارد، پس نیت‌هایش هم خیرخواهانه است» \\
\rowcolor{NeutralLight}
سوگیری تأییدی (\lr{Confirmation Bias}) &
جست‌وجوی شواهد تأییدکننده‌ی باور قبلی &
نادیده‌گرفتن نشانه‌های سوءنیت مداخله‌گر \\
ناامیدی آموخته‌شده (\lr{Learned Helplessness}) &
باور به ناتوانی ذاتی خود در تغییر وضع &
«ما نمی‌توانیم خودمان کاری کنیم، فقط بیگانه قادر است»%
\footnote{\lr{Seligman, Martin E.P. \textit{Helplessness: On Depression, Development, and Death}. Freeman, 1975.}} \\
\rowcolor{NeutralLight}
تفکر سیاه-سفید (\lr{Splitting}) &
تقسیم جهان به کاملاً خوب و کاملاً بد &
«حکومت شر مطلق و مداخله‌گر خیر مطلق است» \\

\end{longtable}

% ─────────────────────────────────────────────────────────────────────────
\section{الگوی ترکیبی پیشنهادی: مدل «شش‌وجهی نجات‌طلبی»}
\label{sec:ch03-hexagonal}
% ─────────────────────────────────────────────────────────────────────────

بر مبنای بررسی الگوهای پنج‌گانه‌ی فوق، الگوی ترکیبی زیر برای تحلیل نظام‌مند هر مورد نجات‌طلبی پیشنهاد می‌شود. این الگو \concept{شش‌وجهی} است زیرا هر مورد را از شش زاویه‌ی متفاوت بررسی می‌کند:

\begin{figure}[htbp]
\centering
\begin{tikzpicture}[
  every node/.style={font=\small},
  hexnode/.style={
    regular polygon, regular polygon sides=6,
    draw=PrimaryMid, fill=BgCool, minimum size=2.8cm,
    font=\small\bfseries, align=center, inner sep=2pt
  },
  outernode/.style={
    draw=PrimaryLight, fill=white, rounded corners=4pt,
    text width=2.8cm, align=center, minimum height=1.2cm,
    font=\scriptsize
  },
  arrow/.style={->, >=Stealth, thick, PrimaryMid!70},
]

% شش‌ضلعی مرکزی
\node[hexnode, fill=AccentGold!25, minimum size=3.5cm] (center)
  {\rl{تحلیل}\\%
   \rl{نجات‌طلبی}};

% شش وجه
\node[outernode, fill=BgWarm] (v1) at (90:4.5cm)
  {\rl{\textbf{۱. ساختاری-سیاسی}}\\%
   \rl{ماهیت رژیم، بحران مشروعیت}};

\node[outernode, fill=BgTeal] (v2) at (30:4.5cm)
  {\rl{\textbf{۲. حقوقی-هنجاری}}\\%
   \rl{مبنای حقوقی، هنجارهای بین‌المللی}};

\node[outernode, fill=BgSuccess] (v3) at (330:4.5cm)
  {\rl{\textbf{۳. اقتصادی}}\\%
   \rl{منابع، وابستگی، هزینه‌ها}};

\node[outernode, fill=BgFail!50] (v4) at (270:4.5cm)
  {\rl{\textbf{۴. نظامی-امنیتی}}\\%
   \rl{توازن قوا، لجستیک، تلفات}};

\node[outernode, fill=BgCool] (v5) at (210:4.5cm)
  {\rl{\textbf{۵. روان‌شناختی-فرهنگی}}\\%
   \rl{هویت، روایت، سوگیری‌ها}};

\node[outernode, fill=AccentPurple!10] (v6) at (150:4.5cm)
  {\rl{\textbf{۶. ژئوپولیتیک}}\\%
   \rl{منافع قدرت‌ها، نظام بین‌الملل}};

% فلش‌ها
\draw[arrow] (center) -- (v1);
\draw[arrow] (center) -- (v2);
\draw[arrow] (center) -- (v3);
\draw[arrow] (center) -- (v4);
\draw[arrow] (center) -- (v5);
\draw[arrow] (center) -- (v6);

% خطوط ارتباطی بین وجوه
\draw[dashed, NeutralMid, thin] (v1) -- (v2);
\draw[dashed, NeutralMid, thin] (v2) -- (v3);
\draw[dashed, NeutralMid, thin] (v3) -- (v4);
\draw[dashed, NeutralMid, thin] (v4) -- (v5);
\draw[dashed, NeutralMid, thin] (v5) -- (v6);
\draw[dashed, NeutralMid, thin] (v6) -- (v1);

\end{tikzpicture}
\caption{مدل شش‌وجهی تحلیل نجات‌طلبی}
\label{fig:ch03-hexagonal}
\end{figure}

\subsection{شرح وجوه شش‌گانه}
\label{subsec:ch03-hex-detail}

\begin{tcolorbox}[casebox, title={وجه ۱: تحلیل ساختاری-سیاسی}]
\textbf{پرسش‌های کلیدی:}
\begin{itemize}[nosep, rightmargin=1em]
  \item ماهیت رژیم حاکم چیست؟ (دموکراتیک، اقتدارگرا، توتالیتر، شکست‌خورده)
  \item شدت سرکوب و نقض حقوق بشر چقدر است؟
  \item آیا کانال‌های داخلی تغییر (انتخابات، اصلاحات، شورش) مسدود شده‌اند؟
  \item ترکیب قومی-مذهبی-طبقاتی دعوت‌کنندگان چیست؟
  \item آیا دعوت‌کنندگان نمایندگی واقعی مردم را دارند؟
\end{itemize}
\textbf{شاخص‌ها:} \lr{Freedom House Index, Polity V, V-Dem, Fragile States Index.}
\end{tcolorbox}

\begin{tcolorbox}[casebox, title={وجه ۲: تحلیل حقوقی-هنجاری}]
\textbf{پرسش‌های کلیدی:}
\begin{itemize}[nosep, rightmargin=1em]
  \item آیا مبنای حقوقی بین‌المللی برای مداخله وجود دارد؟
  \item آیا مجوز شورای امنیت صادر شده؟
  \item آیا دعوت از حکومت مشروع است یا از اپوزیسیون/اقلیت؟
  \item آیا هنجارهای \lr{R2P} قابل اعمال‌اند؟
  \item آیا شرایط «جدایی اصلاحی» فراهم است؟
\end{itemize}
\end{tcolorbox}

\begin{tcolorbox}[casebox, title={وجه ۳: تحلیل اقتصادی}]
\textbf{پرسش‌های کلیدی:}
\begin{itemize}[nosep, rightmargin=1em]
  \item منابع طبیعی و اقتصادی منطقه‌ی هدف چیست؟ (آیا «نفرین منابع» وجود دارد؟)
  \item ساختار اقتصادی وابستگی به مداخله‌گر چگونه خواهد بود؟
  \item هزینه‌های بازسازی چقدر و بر عهده‌ی کیست؟
  \item آیا مداخله‌گر امتیازات اقتصادی مطالبه کرده یا خواهد کرد؟
\end{itemize}
\end{tcolorbox}

\begin{tcolorbox}[casebox, title={وجه ۴: تحلیل نظامی-امنیتی}]
\textbf{پرسش‌های کلیدی:}
\begin{itemize}[nosep, rightmargin=1em]
  \item توازن نظامی میان مداخله‌گر و حکومت هدف چگونه است؟
  \item تلفات مورد انتظار (نظامی و غیرنظامی) چقدر است؟
  \item آیا مداخله به بی‌ثباتی منطقه‌ای منجر خواهد شد؟
  \item استراتژی خروج چیست؟
  \item آیا پایگاه‌های نظامی دائم ایجاد خواهد شد؟
\end{itemize}
\end{tcolorbox}

\begin{tcolorbox}[casebox, title={وجه ۵: تحلیل روان‌شناختی-فرهنگی}]
\textbf{پرسش‌های کلیدی:}
\begin{itemize}[nosep, rightmargin=1em]
  \item کدام سوگیری‌های شناختی در تصمیم دعوت‌کنندگان نقش دارد؟
  \item آیا «ناامیدی آموخته‌شده» بر جامعه حاکم است؟
  \item روایت غالب از «نجات‌دهنده» چیست؟ (آرمانی‌سازی یا واقع‌بینانه)
  \item آیا پیشینه‌ی فرهنگی مشترکی میان دعوت‌کننده و مداخله‌گر وجود دارد؟
  \item چگونه حافظه‌ی تاریخی جامعه بر پذیرش مداخله اثر می‌گذارد؟
\end{itemize}
\end{tcolorbox}

\begin{tcolorbox}[casebox, title={وجه ۶: تحلیل ژئوپولیتیک}]
\textbf{پرسش‌های کلیدی:}
\begin{itemize}[nosep, rightmargin=1em]
  \item منافع ژئوپولیتیک مداخله‌گر چیست؟
  \item واکنش قدرت‌های رقیب چه خواهد بود؟
  \item آیا مداخله در چارچوب رقابت ابرقدرتی صورت می‌گیرد؟
  \item پیامدهای منطقه‌ای مداخله (اثر دومینو، تغییر مرزها) چیست؟
  \item آیا منطقه‌ی هدف اهمیت استراتژیک (تنگه‌ها، مسیرها، منابع) دارد؟
\end{itemize}
\end{tcolorbox}

% ─────────────────────────────────────────────────────────────────────────
\section{الگوی طیفی نتایج مداخله}
\label{sec:ch03-spectrum}
% ─────────────────────────────────────────────────────────────────────────

نتایج مداخلات خارجی را نمی‌توان در دوگانه‌ی ساده‌ی «موفقیت/شکست» خلاصه کرد. الگوی طیفی زیر نشان می‌دهد که نتایج در یک پیوستار قرار دارند:

\begin{figure}[htbp]
\centering
\begin{tikzpicture}[
  every node/.style={font=\small},
]

% خط طیف
\draw[line width=3pt, left color=AccentGreen, right color=AccentRed]
  (0,0) -- (14,0);

% نقاط
\node[success mark] at (1,0) {A};
\node[success mark] at (3.5,0) {B};
\node[mixed mark] at (6,0) {C};
\node[mixed mark] at (8.5,0) {D};
\node[fail mark] at (11,0) {E};
\node[fail mark] at (13,0) {F};

% برچسب‌ها بالا
\node[above=0.8cm, text width=2.5cm, align=center, font=\scriptsize\bfseries,
      color=AccentGreen] at (1,0)
  {\rl{رهایی کامل و}\\%
   \rl{استقلال پایدار}};
\node[above=0.8cm, text width=2.5cm, align=center, font=\scriptsize\bfseries,
      color=AccentGreen!70] at (3.5,0)
  {\rl{رهایی با}\\%
   \rl{وابستگی محدود}};
\node[above=0.8cm, text width=2.5cm, align=center, font=\scriptsize\bfseries,
      color=AccentAmber] at (6,0)
  {\rl{تغییر رژیم اما}\\%
   \rl{بی‌ثباتی}};
\node[above=0.8cm, text width=2.5cm, align=center, font=\scriptsize\bfseries,
      color=AccentAmber!80] at (8.5,0)
  {\rl{رژیم دست‌نشانده}\\%
   \rl{و وابستگی ساختاری}};
\node[above=0.8cm, text width=2.5cm, align=center, font=\scriptsize\bfseries,
      color=AccentRed!70] at (11,0)
  {\rl{هرج‌ومرج و}\\%
   \rl{دولت شکست‌خورده}};
\node[above=0.8cm, text width=2.5cm, align=center, font=\scriptsize\bfseries,
      color=AccentRed] at (13,0)
  {\rl{استعمار یا}\\%
   \rl{الحاق}};

% برچسب‌ها پایین: نمونه‌ها
\node[below=0.8cm, text width=2.5cm, align=center, font=\scriptsize,
      color=NeutralDark] at (1,0)
  {\rl{فرانسه‌ی آزاد}\\%
   \rl{(۱۹۴۴)}};
\node[below=0.8cm, text width=2.5cm, align=center, font=\scriptsize,
      color=NeutralDark] at (3.5,0)
  {\rl{کوزوو}\\%
   \rl{(۱۹۹۹)}};
\node[below=0.8cm, text width=2.5cm, align=center, font=\scriptsize,
      color=NeutralDark] at (6,0)
  {\rl{عراق}\\%
   \rl{(۲۰۰۳)}};
\node[below=0.8cm, text width=2.5cm, align=center, font=\scriptsize,
      color=NeutralDark] at (8.5,0)
  {\rl{افغانستان}\\%
   \rl{(۲۰۰۱)}};
\node[below=0.8cm, text width=2.5cm, align=center, font=\scriptsize,
      color=NeutralDark] at (11,0)
  {\rl{لیبی}\\%
   \rl{(۲۰۱۱)}};
\node[below=0.8cm, text width=2.5cm, align=center, font=\scriptsize,
      color=NeutralDark] at (13,0)
  {\rl{کریمه}\\%
   \rl{(۲۰۱۴)}};

\end{tikzpicture}
\caption{طیف نتایج مداخلات خارجی}
\label{fig:ch03-spectrum}
\end{figure}

% ─────────────────────────────────────────────────────────────────────────
\section{جدول تطبیقی: اِعمال الگوی شش‌وجهی بر موارد کلیدی}
\label{sec:ch03-comparative}
% ─────────────────────────────────────────────────────────────────────────

\begin{landscape}
\begin{table}[htbp]
\centering
\caption{اِعمال الگوی شش‌وجهی بر ده مورد کلیدی}\label{tab:ch03-comparative}
\begin{adjustbox}{max width=\linewidth}
\begin{tabular}{
  >{\raggedleft\arraybackslash\bfseries\small}p{2cm}
  >{\raggedleft\arraybackslash\small}p{3.5cm}
  >{\raggedleft\arraybackslash\small}p{2.5cm}
  >{\raggedleft\arraybackslash\small}p{3.5cm}
  >{\raggedleft\arraybackslash\small}p{3.5cm}
  >{\raggedleft\arraybackslash\small}p{3.5cm}
  >{\raggedleft\arraybackslash\small}p{3.5cm}
  >{\raggedleft\arraybackslash\small}p{2cm}}
\toprule
\rowcolor{PrimaryDeep}
\textcolor{white}{مورد} &
\textcolor{white}{ساختاری} &
\textcolor{white}{حقوقی} &
\textcolor{white}{اقتصادی} &
\textcolor{white}{نظامی} &
\textcolor{white}{روان‌شناختی} &
\textcolor{white}{ژئوپولیتیک} &
\textcolor{white}{نتیجه (طیف)} \\
\midrule

انقلاب شکوهمند (۱۶۸۸) &
\cellcolor{AccentGreen!20}بالا &
\cellcolor{AccentGreen!20}بالا &
\cellcolor{AccentGreen!20}کم‌هزینه &
\cellcolor{AccentGreen!20}بدون خونریزی &
\cellcolor{AccentGreen!20}واقع‌بینانه &
\cellcolor{mixed mark!20}متوسط &
\cellcolor{AccentGreen!20}A \\
\rowcolor{NeutralLight}
فتح عرب ایران (۶۳۶) &
بحران شدید &
نامرتبط &
غارت &
خونین &
شکست روحی &
تغییر تمدنی &
E–F \\
مجاهدین افغان (۱۹۷۹) &
سرکوب شدید &
ضعیف &
هزینه‌ی بالا &
طولانی &
آرزواندیشی &
رقابت ابرقدرتی &
D–E \\
\rowcolor{NeutralLight}
کوزوو (۱۹۹۹) &
بحران شدید &
بدون مجوز &
متوسط &
متناسب &
ناامیدی واقعی &
رقابت منطقه‌ای &
B \\
عراق (۲۰۰۳) &
استبداد شدید &
بسیار ضعیف &
نفرین نفت &
سریع اما باتلاق &
آرزواندیشی &
نفت، اسرائیل &
C–D \\
\rowcolor{NeutralLight}
لیبی (۲۰۱۱) &
استبداد &
قطعنامه ۱۹۷۳ &
نفت &
سریع اما بی‌نقشه &
آرزواندیشی &
اروپا، نفت &
E \\
سوریه (۲۰۱۵) &
جنگ داخلی &
دعوت حکومت &
ویرانی &
طولانی &
یأس &
رقابت جهانی &
D–E \\
\rowcolor{NeutralLight}
کریمه (۲۰۱۴) &
بحران سیاسی &
بسیار ضعیف &
متوسط &
بدون مقاومت &
هویت روسی &
ناتو-روسیه &
F \\
فرانسه‌ی آزاد (۱۹۴۴) &
اشغال &
جنگ جهانی &
مارشال &
خونین اما قاطع &
مقاومت قوی &
جنگ جهانی &
A \\
\rowcolor{NeutralLight}
تیمور شرقی (۱۹۹۹) &
سرکوب شدید &
قطعنامه ۱۲۶۴ &
فقر &
محدود &
واقع‌بینانه &
اندونزی-استرالیا &
A–B \\

\bottomrule
\end{tabular}
\end{adjustbox}
\end{table}
\end{landscape}

% ─────────────────────────────────────────────────────────────────────────
\section{الگوی پیش‌بینی: شرایط لازم و کافی برای موفقیت}
\label{sec:ch03-prediction}
% ─────────────────────────────────────────────────────────────────────────

بر اساس بررسی تطبیقی موارد بالا، می‌توان مجموعه‌ای از \concept{شرایط لازم} و \concept{شرایط کافی} برای «موفقیت نسبی» نجات‌طلبی استخراج کرد:

\begin{tcolorbox}[successbox, title={شرایط لازم (همه باید برقرار باشند)}]
\begin{enumerate}[nosep, label=\textbf{ل\arabic*.}]
  \item ستم و سرکوب مستند و گسترده باشد (نه صرفاً ادعای بخشی از نخبگان).
  \item حداقلی از ظرفیت سازمانی داخلی برای دولت‌سازی پس از مداخله وجود داشته باشد.
  \item مداخله‌گر دارای استراتژی خروج مشخص باشد.
  \item حمایت بین‌المللی گسترده (نه فقط یک قدرت) وجود داشته باشد.
  \item جایگزین سیاسی مشخصی برای رژیم فعلی وجود داشته باشد.
\end{enumerate}
\end{tcolorbox}

\begin{tcolorbox}[defbox, title={شرایط کافی (حضور آنها احتمال موفقیت را به‌شدت افزایش می‌دهد)}]
\begin{enumerate}[nosep, label=\textbf{ک\arabic*.}]
  \item مجوز شورای امنیت یا اجماع منطقه‌ای وجود داشته باشد.
  \item مداخله چندجانبه باشد (نه یک‌جانبه).
  \item جامعه‌ی هدف دارای سنت نهادی (حتی ضعیف) باشد.
  \item مداخله‌گر تعهد مالی بلندمدت به بازسازی داشته باشد.
  \item جمعیت محلی حمایت گسترده از مداخله داشته باشد (نه فقط نخبگان).
  \item هیچ قدرت منطقه‌ای یا جهانی فعالانه مخالف مداخله نباشد.
\end{enumerate}
\end{tcolorbox}

% ─────────────────────────────────────────────────────────────────────────
\section{جمع‌بندی فصل}
\label{sec:ch03-summary}
% ─────────────────────────────────────────────────────────────────────────

این فصل پنج الگوی تحلیلی (واقع‌گرایانه، نهادگرایانه، سازه‌انگارانه، پسااستعماری و روان‌شناسی سیاسی) را بررسی کرد و نقاط قوت و ضعف هرکدام را نشان داد. الگوی ترکیبی \concept{شش‌وجهی} برای تحلیل جامع هر مورد نجات‌طلبی پیشنهاد شد و بر ده مورد تاریخی اِعمال گردید. الگوی \concept{طیفی} نتایج مداخلات نشان داد که دوگانه‌ی ساده‌ی «موفقیت/شکست» گمراه‌کننده است و مجموعه‌ای از شرایط لازم و کافی برای موفقیت نسبی استخراج شد.

در فصل‌های بعد، این ابزارهای تحلیلی بر موارد تاریخی مشخص — ابتدا در ایران و سپس در جهان — اعمال خواهد شد.

\cleardoublepage